\chapter{System Analysis and Design}

\section{Stakeholder Analysis}

Primary stakeholders include visitors, curators, research engineers, and operations
teams. Each requires a different quality guarantee: usability, reliability,
explainability, and maintainability. Visitors prioritize speed and clarity; curators
emphasize factual integrity and update control; engineers prioritize maintainability
and reproducibility; leadership focuses on adoption and measurable engagement. A
system that optimizes only one of these viewpoints tends to fail in operation; hence
requirements engineering here explicitly balances all four.

\section{Functional Requirements}

\begin{enumerate}[label=\textbf{FR-\arabic*.},leftmargin=3em]
  \item \textbf{Responsive detection.} Detect and track artifacts in the live camera
        feed with bounding-box overlays at interactive frame rates.
  \item \textbf{On-device classification.} Classify a framed artifact as one of the
        51 catalogued GEM artifacts entirely on-device, with no per-frame network call.
  \item \textbf{Confidence gating.} A confident detection opens the artifact's
        information; a view that falls below the per-class confidence threshold is
        suppressed, so the app keeps scanning instead of guessing.
  \item \textbf{Artifact information.} Present name, period/dynasty, material, and
        description, with an original$\rightarrow$restored image toggle where available.
  \item \textbf{Bilingual AI guide.} Answer artifact-specific and general questions in
        the user's language, grounded in curated museum data.
  \item \textbf{Bilingual UI.} Support full English/Arabic switching including
        right-to-left layout, with the chosen language persisted across sessions.
  \item \textbf{Text-to-speech.} Optionally read AI replies aloud in the active language.
  \item \textbf{Photo capture.} Capture and save camera photos to the device gallery.
  \item \textbf{Museum information and location.} Provide a museum-information screen
        (hours, location, ticketing, map) with location gating for on-site versus
        off-site visitors.
\end{enumerate}

\section{Non-Functional Requirements}

\begin{enumerate}[label=\textbf{NFR-\arabic*.},leftmargin=3.4em]
  \item \textbf{Latency:} real-time on-device detection (target $\geq 2.5$\,FPS
        on a mid-tier Android phone) and AI answers within a few seconds.
  \item \textbf{Footprint:} compact TensorFlow Lite model ($\leq 8$\,MB; actual
        5.97\,MB), installable on low-storage devices.
  \item \textbf{Robustness:} tolerate clutter, overlap, oblique angles, and
        fine-grained confusion among near-identical artifacts.
  \item \textbf{Privacy/edge:} perform recognition on-device with no per-frame
        upload; serve the AI backend locally, avoiding any third-party LLM API.
  \item \textbf{Accessibility:} bilingual content plus text-to-speech audio.
  \item \textbf{Maintainability/scalability:} modular, independently replaceable
        services; the knowledge base is updatable without an app release.
\end{enumerate}

\section{Software and Hardware Requirements}

\begin{table}[htbp]
    \centering
    \caption{Mobile Client (on-device AI) Software \& Hardware Requirements}
    \label{tab:req-mobile}
    \renewcommand{\arraystretch}{1.3}
    \begin{tabular}{>{\bfseries}p{3.5cm} p{8.5cm}}
        \toprule
        \rowcolor{gemsand}
        \textbf{Category} & \textbf{Requirement} \\
        \midrule
        OS
            & Android 8.0+ / IOS 13.0+ \\
        RAM
            & 298MB \\
        Storage
            & 45MB \\
        Processor
            & Samsung Exynos 1580 / A14 Bionic \\
        Network
            & Required only for IR backend calls \\
        \bottomrule
    \end{tabular}
\end{table}

\begin{table}[htbp]
    \centering
    \caption{Development Environment Software \& Requirements}
    \label{tab:req-devenv}
    \renewcommand{\arraystretch}{1.3}
    \begin{tabular}{>{\bfseries}p{5.5cm} p{6.5cm}}
        \toprule
        \rowcolor{gemsand}
        \textbf{Category} & \textbf{Requirement} \\
        \midrule
        IDE
            & Android Studio / Visual Studio Code / Xcode \\
        ML Runtime
            & Python 3.12 / Flutter 3.41 \\
        Framework
            & FastAPI / Ultralytics~8.4.45 / PyTorch TensorFlow~Lite (FP16) / LlamaIndex / LangChain \\
        Detector Training
            & NVIDIA GTX 1650 Ti (4GB) \\
        RAG Experiments
            &  NVIDIA RTX 3050 Laptop (4GB) \\
        Production Serving
            &  RunPod NVIDIA RTXA5000 (24GB) \\
        \bottomrule
    \end{tabular}
\end{table}

% \begin{table}[htbp]
%     \centering
%     \caption{Backend (IR Server) Software \& Hardware Requirements}
%     \label{tab:req-backend}
%     \renewcommand{\arraystretch}{1.3}
%     \begin{tabular}{>{\bfseries}p{3.5cm} p{8.5cm}}
%         \toprule
%         \rowcolor{tableheader}
%         \textbf{Category} & \textbf{Requirement} \\
%         \midrule
%         Server OS
%             & Ubuntu X.X LTS / Windows Se \\
%         Runtime Environment
%             & Python 3.X / Node.js X.X (specify) \\
%         IR Framework
%             & Elasticsearch X.X / Whoosh X.X / BM25 library (specify) \\
%         Network
%             & HTTPS endpoint; minimum bandwidth sufficient to meet
%               query-response latency NFR \\
%         Hardware (Self-hosted)
%             & Minimum X-core CPU, N\,GB RAM, M\,GB SSD; or cloud
%               equivalent (e.g., AWS t3.medium) \\
%         \bottomrule
%     \end{tabular}
% \end{table}

\FloatBarrier
\section{Data Requirements}

\subsection{Data Sources and Acquisition}
Training data consists of roughly 30{,}000 frames extracted (at about one frame per
second) from approximately 50 in-gallery videos of GEM artifacts captured during
the museum's opening hours, spanning all 51 classes. This source reflects the real
deployment domain far better than generic web imagery.

\subsection{Volumetric and Diversity Requirements}
The dataset must cover class diversity, view diversity, and clutter diversity, with
balanced per-class representation. Per-class frame counts range from a few hundred to
just under a thousand, which is reasonably balanced for the task.

\subsection{Annotation and Labeling Standards}
Annotation consistency, class-map governance, and \emph{leakage-aware} split policies
are required. Because frames are extracted from continuous video, adjacent frames are
near-duplicates; splits must therefore be performed \emph{chronologically within each
class/video group} rather than by random shuffle, to avoid leaking near-identical
frames across train, validation, and test (see Chapter~5).

\section{App Persona}

To keep design decisions anchored to a concrete user rather than an abstract average,
we model a representative visitor (Figure~\ref{fig:persona}). Her goals and frustrations
motivate the system's emphasis on speed, offline robustness, and a low-friction
scan-to-answer path.

\begin{figure}[htbp]
\centering
\fbox{\begin{minipage}{0.9\textwidth}\small
\vspace{1mm}
{\large\bfseries Persona: Sarah the Curious Visitor}\\[1mm]
\textit{32 years old, international tourist, tech-savvy but time-constrained.}\\[2mm]
\begin{minipage}[t]{0.46\textwidth}
\textbf{Goals:}
\begin{itemize}[leftmargin=1.2em,nosep]
  \item Quickly identify artifacts.
  \item Learn historical context.
  \item Interactive discovery.
\end{itemize}
\end{minipage}\hfill
\begin{minipage}[t]{0.46\textwidth}
\textbf{Frustrations:}
\begin{itemize}[leftmargin=1.2em,nosep]
  \item Weak Wi-Fi in halls.
  \item Static museum text.
  \item High app latency.
\end{itemize}
\end{minipage}\\[2mm]
\textit{``I want to point my phone at something and instantly hear its story, like a
personal digital curator.''}
\vspace{1mm}
\end{minipage}}
\caption{User persona modeling for the \TimeLens{} experience design.}
\label{fig:persona}
\end{figure}

\section{Use Case Diagram}

The visitor is the primary actor, and the system exposes three core use cases that map
directly to the functional requirements (Figure~\ref{fig:usecase}): scanning to
identify an artifact (FR-01--FR-03), viewing its information (FR-04), and asking the
bilingual AI guide a question (FR-05). All three are reached from a single camera-first
screen, keeping the interaction model deliberately small.

\begin{figure}[htbp]
\centering
\begin{tikzpicture}[node distance=6mm and 28mm]
  % actor
  \node (head) at (0,0) [circle,draw,minimum size=5mm] {};
  \draw (0,-2.5mm)--(0,-12mm);            % body
  \draw (-5mm,-6mm)--(5mm,-6mm);          % arms
  \draw (0,-12mm)--(-4mm,-20mm);          % legs
  \draw (0,-12mm)--(4mm,-20mm);
  \node (visitor) at (0,-2.5) {\small Visitor};
  % use cases
  \node[draw,ellipse,align=center,minimum width=34mm,font=\small] (scan) at (4.2,1.2) {Scan / Identify\\ Artifact};
  \node[draw,ellipse,align=center,minimum width=34mm,font=\small] (info) at (4.2,-0.3) {View Artifact\\ Information};
  \node[draw,ellipse,align=center,minimum width=34mm,font=\small] (chat) at (4.2,-1.8) {Ask the AI Guide};
  % system boundary
  \begin{scope}[on background layer]
    \node[draw,dashed,rounded corners,fit=(scan)(info)(chat),inner sep=5mm,label=above:{\small TimeLens Platform}] {};
  \end{scope}
  % associations
  \draw (head) -- (scan);
  \draw (head) -- (info);
  \draw (head) -- (chat);
\end{tikzpicture}
\caption{Primary use cases for visitor interaction.}
\label{fig:usecase}
\end{figure}

\FloatBarrier

\section{Interaction and Data Modeling}

System interactions are organized by actor role. Data modeling focuses on the
alignment of user sessions, artifact scans, and bilingual content, ensuring
consistency between the on-device layer and the knowledge backend.

\section{Requirements Traceability Matrix}

Table~\ref{tab:traceability} maps a representative subset of requirements to the
components that satisfy them and the method used to validate each.

\begin{table}[htbp]
\centering\small
\caption{Requirement-to-component mapping.}
\label{tab:traceability}
\begin{tabularx}{\textwidth}{@{}l>{\raggedright\arraybackslash}X>{\raggedright\arraybackslash}X>{\raggedright\arraybackslash}p{2.3cm}@{}}
\toprule
\textbf{Req.} & \textbf{Description} & \textbf{Component} & \textbf{Validation}\\
\midrule
FR-01 & Live detection $+$ tracking overlay & ARService $+$ TFLiteService (YOLOv8n) & Functional test\\
FR-02 & On-device 51-class classification & TFLiteService (end-to-end NMS) & Per-class mAP eval\\
FR-03 & Confident $\rightarrow$ info; sub-threshold $\rightarrow$ suppressed & Home / Scanning logic & End-to-end scenario\\
FR-05 & Bilingual grounded AI guide & RAG \texttt{/ask} (ChromaDB $+$ Gemma) $+$ ChatDrawer & API integration\\
FR-06 & EN/AR switch $+$ RTL $+$ persistence & LocaleController $+$ l10n & UI test (both locales)\\
FR-07 & Text-to-speech replies & \texttt{flutter\_tts} (ChatDrawer) & Functional test\\
NFR-2 & Compact model $\leq 8$\,MB & YOLOv8n TFLite FP16 (5.97\,MB) & Asset-size check\\
NFR-4 & On-device privacy / local backend & On-device inference $+$ local Ollama & Architecture review\\
\bottomrule
\end{tabularx}
\end{table}

\FloatBarrier
\section{High-Level Architecture}

The architecture is intentionally split into two layers with a clear boundary. On the
phone, an \emph{on-device perception layer} turns each camera frame into a tracked,
classified artifact. A \emph{knowledge layer} then grounds any follow-up question in
curated museum data before a language model phrases the answer. This framing avoids a
common failure mode in heritage systems, where a single weak prediction directly
triggers a confident generated answer.

At the system level, \TimeLens{} separates concerns across five layers: experience,
on-device detection, knowledge and retrieval, service, and data. Each layer can evolve
independently behind a stable contract: the detector can be retrained, the knowledge
base updated, or the language model swapped, without collapsing the end-to-end journey.

\section{High-Level System Block Diagram}

Figure~\ref{fig:block} shows how the two layers connect at runtime. Detection runs
entirely on the phone; only a chat question crosses the network, reaching the FastAPI
service that retrieves from ChromaDB and prompts the locally served Gemma model.

\begin{figure}[htbp]
\centering
\begin{tikzpicture}[node distance=12mm and 20mm]
  \node[block] (app) {Mobile App (Flutter)\\[1pt]\scriptsize on-device YOLOv8n (TFLite)};
  \node[block,right=of app] (api) {FastAPI\\ RAG Service};
  \node[block,right=of api] (gemma) {Ollama\\ Gemma 4 E2B (Q4)};
  \node[store,below=of api] (chroma) {ChromaDB\\\scriptsize EN/AR collections};
  \draw[flow] ([yshift=2mm]app.east) -- node[above,font=\scriptsize]{HTTPS / ngrok} ([yshift=2mm]api.west);
  \draw[flow] ([yshift=-2mm]api.west) -- node[below,font=\scriptsize]{answer} ([yshift=-2mm]app.east);
  \draw[flow] (api) -- node[above,font=\scriptsize]{prompt} (gemma);
  \draw[flow] (api) -- node[right,font=\scriptsize]{retrieve} (chroma);
\end{tikzpicture}
\caption{System block diagram: detection runs on-device; only chat questions reach the
FastAPI RAG service, which retrieves from ChromaDB and prompts a locally served Gemma
model.}
\label{fig:block}
\end{figure}

\FloatBarrier

\section{Layered System Architecture}

Viewed as layers (Figure~\ref{fig:layers}), the separation of concerns becomes
explicit. The experience layer fans out into two independent paths: an
\emph{on-device path}, where detection resolves artifact identity and metadata
with no network dependency, and a \emph{service path}, where a chat question
crosses HTTPS/ngrok once to the FastAPI endpoints fronting the knowledge layer.
Each boundary is a stable contract, and both paths terminate in the same
canonical dataset, so \emph{grounding}---constraining generation to retrieved
museum data---holds by construction.

\begin{figure}[htbp]
\centering
\begin{tikzpicture}[every node/.style={font=\small}]
  % column layers
  \node[block,text width=50mm,minimum height=11mm] (det) at (0,0)
    {On-Device Detection Layer\\[1pt]\scriptsize YOLOv8n, TFLite, end-to-end NMS};
  \node[block,text width=54mm,minimum height=11mm] (svc) at (7.4,0)
    {Chatbot Service Layer\\[1pt]\scriptsize FastAPI: \texttt{/ask}, \texttt{/artifact}, \texttt{/health}};
  \node[edgenode,text width=54mm,minimum height=11mm] (rag) at (7.4,-2.6)
    {Knowledge \& RAG Layer\\[1pt]\scriptsize ChromaDB retrieval $+$ Gemma 4 E2B (Q4) via Ollama};
  % spanning layers
  \node[block,text width=127mm,minimum height=11mm] (exp) at (3.7,2.3)
    {Experience Layer\\[1pt]\scriptsize Flutter UI, camera overlay, bilingual EN/AR, voice playback};
  \node[store,text width=127mm,minimum height=11mm] (data) at (3.7,-5.8)
    {Data Layer: one canonical 108-record bilingual KB\\[1pt]\scriptsize
     bundled JSON (on-device) $=$ ChromaDB \texttt{gem\_en}/\texttt{gem\_ar} (backend);
     preferences on-device};
  % backend boundary
  \begin{scope}[on background layer]
    \node[draw=gemobsidian!45,dashed,rounded corners=4pt,inner sep=3mm,
          fit=(svc)(rag)] (bk) {};
  \end{scope}
  \node[font=\scriptsize\itshape,text=gemobsidian!75,anchor=north west,align=left]
    at ([yshift=-0.5mm]bk.south west) {backend\\(RunPod A5000)};
  \node[font=\scriptsize\itshape,text=gemobsidian!75,anchor=north west]
    at ([yshift=-0.5mm]det.south west) {on-device};
  % layer contracts: requests flow down, responses flow up
  \draw[flow] ($(exp.south -| det)+(-2.5mm,0)$) -- node[left,font=\scriptsize]{frames}
              ($(det.north)+(-2.5mm,0)$);
  \draw[flow] ($(det.north)+(2.5mm,0)$) -- node[right,font=\scriptsize]{boxes}
              ($(exp.south -| det)+(2.5mm,0)$);
  \draw[flow] ($(exp.south -| svc)+(-2.5mm,0)$) -- node[left,align=right,font=\scriptsize]
              {question\\ \textit{HTTPS / ngrok}} ($(svc.north)+(-2.5mm,0)$);
  \draw[flow] ($(svc.north)+(2.5mm,0)$) -- node[right,font=\scriptsize]{answer}
              ($(exp.south -| svc)+(2.5mm,0)$);
  \draw[flow] ($(det.south)+(-2.5mm,0)$) -- node[left,font=\scriptsize]{label}
              ($(data.north -| det)+(-2.5mm,0)$);
  \draw[flow] ($(data.north -| det)+(2.5mm,0)$) -- node[right,font=\scriptsize]{metadata}
              ($(det.south)+(2.5mm,0)$);
  \draw[flow] ($(svc.south)+(-2.5mm,0)$) -- node[left,font=\scriptsize]
              {question $+$ \texttt{artifact\_id}} ($(rag.north)+(-2.5mm,0)$);
  \draw[flow] ($(rag.north)+(2.5mm,0)$) -- node[right,font=\scriptsize]{grounded answer}
              ($(svc.south)+(2.5mm,0)$);
  \draw[flow] ($(rag.south)+(-2.5mm,0)$) -- node[left,font=\scriptsize,pos=0.75]{query}
              ($(data.north -| rag)+(-2.5mm,0)$);
  \draw[flow] ($(data.north -| rag)+(2.5mm,0)$) -- node[right,font=\scriptsize,pos=0.25]{records EN/AR}
              ($(rag.south)+(2.5mm,0)$);
  % grounding guarantee on the generative path
  \node[rotate=-90,anchor=center,font=\footnotesize\bfseries]
    at (11.3,-2.9) {Grounding};
\end{tikzpicture}
\caption{Layered architecture. Downward arrows carry requests, upward arrows the
responses. The experience layer fans out into two independent
paths: detection and metadata lookup resolve entirely on-device, while a chat
question crosses the network once (HTTPS/ngrok) to the FastAPI service fronting
the RAG engine. Both paths draw on the same canonical 108-record bilingual
dataset---bundled as JSON on the device, indexed as ChromaDB collections on the
backend---and ``grounding'' constrains generation to that retrieved museum data.}
\label{fig:layers}
\end{figure}

\FloatBarrier

\subsection{Class Design}

Figure~\ref{fig:classdesign} sketches the core classes behind this flow. The mobile
client drives an on-device \texttt{TFLiteDetector}, resolves artifact metadata locally,
and calls the bilingual \texttt{RagService} only for chat---there is no backend ensemble
verifier in the deployed path.

\begin{figure}[htbp]
\centering
\begin{tikzpicture}[node distance=8mm and 12mm,cls/.style={draw,fill=gemblue!5,font=\scriptsize,align=left}]
  \node[cls] (client) {\begin{tabular}{@{}l@{}}\textbf{MobileClient}\\\hline currentFrame: Image\\\hline scan()\\ openInfo(id)\\ openChat(id)\end{tabular}};
  \node[cls,right=of client] (det) {\begin{tabular}{@{}l@{}}\textbf{TFLiteDetector}\\\hline modelPath; inputSize=640\\ confidenceThreshold\\\hline parseEndToEndNMS()\\ classify(): Detection\end{tabular}};
  \node[cls,below=of det] (meta) {\begin{tabular}{@{}l@{}}\textbf{ArtifactMetadataService}\\\hline dataset: JSON (51)\\\hline lookup(label): Artifact\end{tabular}};
  \node[cls,right=of det] (rag) {\begin{tabular}{@{}l@{}}\textbf{RagService}\\\hline chroma\_en, chroma\_ar\\ ollama: Gemma4E2B\\\hline detectLanguage()\\ retrieve(); generateAnswer()\end{tabular}};
  \draw[flow] (client) -- (det);
  \draw[flow] (det) -- (meta);
  \draw[flow] (det) -- (rag);
\end{tikzpicture}
\caption{Core class design: on-device detection (\texttt{TFLiteDetector}), local metadata
lookup, and the bilingual RAG service. There is no backend ensemble verifier.}
\label{fig:classdesign}
\end{figure}

\FloatBarrier
\subsection{Scan-to-Answer Sequence}

Figure~\ref{fig:sequence} traces one interaction from end to end. Identity is decided
on-device the moment the visitor frames an artifact; the backend is contacted only when
a question is asked, and only to retrieve context and phrase the answer.

\begin{figure}[htbp]
\centering
\begin{tikzpicture}[font=\scriptsize]
  \def\ytop{0} \def\ybot{-6.2}
  \foreach \x/\name in {0/Visitor, 3/Mobile App, 6.4/FastAPI, 9.6/{ChromaDB + Gemma}}{
    \node[draw,fill=gemsand,minimum width=18mm] (h\x) at (\x,\ytop) {\name};
    \draw[dashed] (\x,\ytop-0.4) -- (\x,\ybot);
  }
  \draw[flow] (0,-0.9) -- node[above]{point camera / ask} (3,-0.9);
  \draw[flow] (3,-1.5) -- ++(0.9,0) |- (3,-1.9) node[right,xshift=10mm]{on-device YOLOv8n detect + classify};
  \draw[flow] (3,-2.5) -- node[above]{POST /ask \{question, artifact\_id\}} (6.4,-2.5);
  \draw[flow] (6.4,-3.1) -- node[above]{retrieve (language-matched)} (9.6,-3.1);
  \draw[flow] (9.6,-3.7) -- node[above]{context + grounded answer} (6.4,-3.7);
  \draw[flow] (6.4,-4.3) -- node[above]{\{answer, language, category\}} (3,-4.3);
  \draw[flow] (3,-4.9) -- node[above]{display + optional TTS} (0,-4.9);
\end{tikzpicture}
\caption{End-to-end scan-to-answer sequence. Identity is decided on-device; the backend
only grounds and phrases the answer.}
\label{fig:sequence}
\end{figure}

\section{System Interaction Timeline}

The interaction is governed by an on-device confidence gate followed by
retrieval-based factual grounding. Every step, from camera capture to the final
answer, is designed so that the language model is invoked only after an artifact
identity is established and only with retrieved, verified context.

\section{Operational Pipeline}

The deployed behavior is supported by a small set of operational modules on the
knowledge side. Rather than a fabricated ``self-healing'' pipeline, these are the real
scripts that build and serve the bilingual knowledge base
(Table~\ref{tab:ragmodules}).

\begin{table}[htbp]
\centering
\caption{Operational modules of the RAG knowledge service.}
\label{tab:ragmodules}
\begin{tabularx}{\textwidth}{@{}llX@{}}
\toprule
\textbf{Stage} & \textbf{Module} & \textbf{Responsibility}\\
\midrule
Dataset cleaning      & \texttt{clean\_dataset.py} & Normalize fields (e.g., strip a stray ``pigment'' prefix from a location).\\
Loading / chunking    & \texttt{load\_dataset.py}  & Split each record into identity / description / significance chunks.\\
Indexing              & \texttt{index\_dataset.py} & Build the language-separated EN/AR ChromaDB collections (215 chunks).\\
Model serving         & \texttt{load\_model.py}    & Load Gemma 4 E2B via Ollama (cached; \texttt{keep\_alive}\,$=-1$).\\
Retrieval + generation& \texttt{Rag\_engine.py}    & Language detect, query classify, retrieve, build prompt, generate.\\
API                   & \texttt{api.py}            & FastAPI \texttt{/ask}, \texttt{/ask/stream}, \texttt{/artifact/\{id\}}, \texttt{/health}.\\
\bottomrule
\end{tabularx}
\end{table}

\section{Inference Pipeline Architecture}

On-device recognition follows a fixed, deterministic sequence
(Figure~\ref{fig:infpipeline}). Each camera frame arrives as sensor-rotated YUV420
data, which is converted to RGB, rotated $90^{\circ}$ to match the Android sensor
orientation, resized to $640\times640$, and normalized before inference. The YOLOv8n
model is exported with non-maximum suppression baked in, so a single forward pass
returns a fixed-size \texttt{[1,300,6]} tensor of \texttt{(xyxy, score, class)} rows
rather than raw anchors, removing any post-processing decode step on the device. A
confidence and geometry filter (global threshold $0.75$, raised to $0.90$ for the
visually near-identical Statue of Siptah) then keeps only trustworthy detections, which
are tracked and drawn as overlays; tapping one opens the artifact's information and the
RAG chat.

\begin{figure}[htbp]
\centering
\begin{tikzpicture}[node distance=7mm,every node/.style={font=\small}]
  \node[block,minimum width=62mm] (a) {Camera frame (YUV420, sensor-rotated)};
  \node[block,minimum width=62mm,below=of a] (b) {YUV $\rightarrow$ RGB conversion};
  \node[block,minimum width=62mm,below=of b] (c) {Rotate $90^{\circ}$ (Android) $\rightarrow$ resize $640$ $\rightarrow$ normalize};
  \node[edgenode,minimum width=62mm,below=of c] (d) {YOLOv8n (end-to-end NMS) $\rightarrow$ \texttt{[1,300,6]}};
  \node[block,minimum width=62mm,below=of d] (e) {Filter: conf $\geq 0.75$ (Siptah $0.90$) $+$ geometry sanity};
  \node[store,minimum width=62mm,below=of e] (f) {Tracked label + box overlay $\rightarrow$ (tap) info + RAG chat};
  \foreach \x/\y in {a/b,b/c,c/d,d/e,e/f}{\draw[flow] (\x) -- (\y);}
\end{tikzpicture}
\caption{On-device inference pipeline from raw camera frame to a tracked, classified
artifact. Non-maximum suppression is baked into the exported model.}
\label{fig:infpipeline}
\end{figure}


\section{Data Flow Architecture}

The information lifecycle is a progressive journey from a raw scan to a grounded
answer: detect on-device, look up local metadata, and---only when the visitor
asks---retrieve bilingual context and generate a response. No camera frames leave the
device, and no per-frame data is logged remotely.

\section{Technology Stack Selection}

\begin{itemize}[leftmargin=1.5em]
  \item \textbf{Mobile app:} Flutter with Provider state management.
  \item \textbf{On-device detection:} YOLOv8n (Ultralytics) exported to TensorFlow Lite
        (FP16) with end-to-end NMS.
  \item \textbf{RAG backend:} FastAPI $+$ LlamaIndex $+$ LangChain; Ollama serving
        Gemma 4 E2B (Q4\_K\_M); ChromaDB vector store; multilingual MiniLM embeddings.
  \item \textbf{Detector training:} Ultralytics 8.4.45 / PyTorch; YOLO-World for
        auto-annotation; Roboflow for hand annotation.
  \item \textbf{Serving / infrastructure:} RunPod-hosted NVIDIA RTX A5000
        (24\,GB VRAM) exposed via an ngrok tunnel.
\end{itemize}

\section{Gemma Integration and Grounding}

Gemma is intentionally downstream of identity. The system never asks Gemma to decide
which artifact is in view; identity is established by the calibrated on-device
detector. Gemma is asked only to narrate \emph{retrieved} museum context after an
artifact is known. This grounding contract keeps the language layer aligned with
museum-safe behavior: if the relevant facts are not in the retrieved context, the
model is instructed to say so rather than invent them.

\section{Production Deployment Strategy}

The deployment defines clear boundaries between the mobile edge (on-device detection
and UI) and the knowledge backend (FastAPI RAG on a RunPod RTX A5000, exposed via ngrok).
This separation keeps the visitor-facing experience responsive and lets the knowledge
base and model evolve without an app release.

\section{App UI and UX Design Integration}

UI and UX were treated as part of system reliability. The app is designed so that user
trust and cognitive load remain stable while model confidence changes in the
background. Instead of exposing raw uncertainty, the navigation guides visitors through
understandable modes: scan, details, contextual learning, and optional conversation.

\section{UI/UX Design Philosophy}

The design philosophy centers on visitor trust. By mapping the operational steps from
visual capture to grounded answer, the system creates a seamless educational journey,
and the hierarchical arrangement of screens lets visitors move naturally between
discovery, scanning, and interactive chat.
